
\section{Preliminary}

This section defines the data format, task views, and unified notation used by DHC. 
\subsection{Dynamic Heterogeneous Graph}

\textbf{Definition 1. Unified dynamic heterogeneous graph.}
We represent a dynamic heterogeneous graph as a chronological event stream
\begin{equation}
\mathcal{S}=\{o_i=(u_i,r_i,v_i,t_i,\mathbf{e}_i,y_i)\}_{i=1}^{N}, \quad t_i \leq t_{i+1},
\label{eq:event_stream}
\end{equation}
where $u_i$ and $v_i$ are the source and destination nodes connected with relation type $r_i\in\mathcal{R}$ which interact at time $t_i$, 
$\mathbf{e}_i$ denotes edge features, and $y_i\in\{0,1\}$ indicates whether the event is a supervised target edge. 
Each node $u\in\mathcal{V}$ has a node type $\tau_u\in\mathcal{U}$ and a text-derived feature vector $\mathbf{x}_u\in\mathbb{R}^{d_x}$. 
The event stream is sorted by time strictly. This definition captures three types of information: typed entities, typed temporal interactions, and text-derived node attributes.

\begin{table}[t]
\centering
\caption{Main notation used in DHC.}
\label{tab:notation}
\scriptsize
\setlength{\tabcolsep}{4pt}
\begin{tabular}{ll}
\toprule
Symbol & Description \\
\midrule
$\mathcal{V},\mathcal{R},\mathcal{U}$ & Node, relation-type, and node-type sets \\
$o_i=(u_i,r_i,v_i,t_i,\mathbf{e}_i,y_i)$ & Temporal event with edge feature and target-edge indicator \\
$\mathbf{x}_u,\tau_u$ & Text-derived node feature and node type of $u$ \\
$\mathbf{M}_u(t)$ & Node memory state of node $u$ after events up to time $t$ \\
$\mathbf{R}_r(t)$ & Relation memory state of relation $r$ after events up to time $t$ \\
$\Phi(\Delta t)$ & Sinusoidal encoding of an elapsed time gap \\
$\rho_r(t)$ & Effective relation representation at time $t$ \\
$\mathcal{N}_{L}(u,t)$ & Bounded recent-neighbor window of node $u$ before time $t$ \\
$\mathbf{E}_u(t)$ & \rev{Sequential event-memory state of node $u$ after events up to time $t$} \\
$\mathbf{q}^{E}_{u\leftarrow v}(t)$ & \rev{Directional message used by the sequential event-memory branch} \\
$h_u^{C}(t), h_u^{E}(t)$ & \rev{Context-branch and sequential event-memory representations} \\
$s_{\theta}(u,r,v,t)$ & Logit score for candidate temporal link $(u,r,v,t)$ \\
$\mathcal{N}_i^{-}$ & Negative destination set sampled for event $o_i$ \\
\bottomrule
\end{tabular}
\end{table}

\subsection{Temporal Edge-Task Views}

\textbf{Definition 2. Temporal heterogeneous edge-task view.}
Given the event stream $\mathcal{S}$, DHC separates events into context edges and target edges. 
Context edges provide temporal and heterogeneous background information for state updates, 
while target edges provide supervised labels for link prediction. 
Formally, events with $y_i=1$ are used in the link-prediction loss, whereas events with $y_i=0$ are used only to update dynamic states.

We consider two task views. In the \emph{transductive} view, training, validation, and test edges are split chronologically, and their node sets may overlap. 
In the \emph{inductive} view, the split also withholds a set of source-side or target-side nodes. Target edges involving these unseen nodes are assigned to the test set. 
Training and validation context edges that touch unseen nodes are removed. This prevents indirect leakage through auxiliary relations and makes the inductive task stricter than ordinary chronological extrapolation.

\subsection{Temporal Link Prediction}

\textbf{Definition 3. Temporal link prediction.}
Given a query $(u,r,?,t)$ and the historical event set
\begin{equation}
\mathcal{H}_t=\{o_i\in\mathcal{S}\mid t_i<t\},
\end{equation}
the goal is to rank the true destination node $v$ above negative candidates. 
DHC computes a temporal link score
\begin{equation}
s_{\theta}(u,r,v,t)=f_{\theta}(u,r,v,t,\mathcal{H}_t),
\end{equation}
where $f_{\theta}$ denotes the DHC scoring function conditioned on dynamic node states, relation states, and recent temporal context built from $\mathcal{H}_t$.

For evaluation, each positive target edge is compared with sampled negative destinations $\mathcal{N}_i^{-}$. 
The rank of the true destination is:
\begin{equation}
\operatorname{rank}_i
=
1+
\sum_{v^{-}\in\mathcal{N}_i^{-}}
\mathbb{I}
\left[
s_{\theta}(u_i,r_i,v^{-},t_i)
\geq
s_{\theta}(u_i,r_i,v_i,t_i)
\right],
\end{equation}
where $\mathbb{I}[\cdot]$ is the indicator function. The mean reciprocal rank(MRR) is
\begin{equation}
\mathrm{MRR}
=
\frac{1}{|\mathcal{Q}|}
\sum_{i\in\mathcal{Q}}
\frac{1}{\operatorname{rank}_i}.
\end{equation}
where $\mathcal{Q}$ is the evaluated positive target-edge set. \rev{Table~\ref{tab:notation} summarizes the symbols used in this paper; in particular, the edge-type augmentation flag and optional relation loss weights are introduced in Algorithm~\ref{alg:dhc_forward} and Eq.~\eqref{eq:link_loss}.}


